\documentclass[french,11pt]{book}
\usepackage{teach}
\usepackage[french]{babel}
\usepackage{amsmath,amssymb}
\usepackage{array}
\newcommand{\R}{\mathbb{R}}
\begin{document}
\begin{center}
{\Large \textbf{Devoir surveillé -- Loi binomiale}}\\[2mm]
Terminale technologique
\end{center}

\section*{Devoir surveillé}
\begin{enumerate}
  \item Une compagnie affirme que 92\% de ses colis arrivent à l'heure. On observe 15 colis. Justifier que l'on peut utiliser une loi binomiale et préciser ses paramètres.
  \item Calculer la probabilité que les 15 colis arrivent à l'heure.
  \item Calculer la probabilité qu'au moins 13 colis arrivent à l'heure.
  \item Déterminer l'espérance et interpréter le résultat.
\end{enumerate}

\end{document}
